\documentclass[letterpaper,10pt]{article}
\usepackage[left=2cm,top=3cm,right=2cm,bottom=4cm]{geometry} 
\usepackage[utf8]{inputenc}
\usepackage[activeacute,spanish]{babel}
\usepackage{multirow}	% Para poder unir filas en las tablas
\usepackage{colortbl}	% Para colorear tablas
\usepackage{latexsym,amssymb,oldlfont,amsmath}
\usepackage{graphicx,graphics}
\usepackage{multicol}
\usepackage{color}
\usepackage{listings}
\usepackage{lastpage}
\usepackage{enumitem}



%\usepackage{fontspec}
%\setmainfont{Times New Roman}
\newcommand{\R}{{\mathbb R}}
\newcommand{\si}{{\mathcal S}^n}
\renewcommand{\theequation}{\thesection.\arabic{equation}} % Cambia formato salida numeracion equation
\makeatletter\@addtoreset {equation}{section} \makeatother % Tambien se usa
\renewcommand{\baselinestretch}{1.3}
\usepackage{amssymb}
\usepackage{hyperref}
\definecolor{midnightblue}{rgb}{0.1,0.1,0.44}


\newcommand{\dis}{\displaystyle}
\renewcommand{\theenumii}{\arabic{enumii}}

\def\Z{\mathbb{Z}}
\def\N{\mathbb{N}}
\def\R{\mathbb{R}}
\def\M{\mathcal{M}_n(\mathbb{R})}
\def\sin{\operatorname{sen}}
\def\e{\operatorname{e}}
\DeclareMathOperator{\arcsec}{arcsec}
\DeclareMathOperator{\arccot}{arccot}
\DeclareMathOperator{\arccsc}{arccsc}
\providecommand{\abs}[1]{\lvert#1\rvert}


\usepackage{etoolbox,fancyhdr,xcolor}
\fancyhf{} % sets both header and footer to nothing
\renewcommand{\headrulewidth}{0pt}
%\renewcommand{\footrulewidth}{0.5pt}
\pagestyle{fancy}
\headheight=60pt %para cambiar el tamaño del encabezadox
\fancyhead[L] %la "L" indica a la izquierda
{
\begin{minipage}{6cm}
\includegraphics[width=0.9\textwidth]{logo_udp.eps}
\end{minipage}
}

%\fancyhead[C]
%{
%\textbf{Cálculo II - Guía 1}
%}


%\fancyhead[R] %la "L" indica a la izquierda
%{

%}
%\begin{minipage}{4cm}
%\small %tamaño de letra pequeño
%{
%\textsf
%{
%Universidad de Chile\\
%...}}
%\end{minipage}}



%\fancyfoot[C]{}
%{
%\begin{minipage}{10cm}
%\color{midnightblue}
%\begin{center}Casa Matriz: Monjitas 527 Piso 15 - Santiago de Chile\\
%Fono (02)24367255 \end{center}
%\end{minipage}
%}
\fancyfoot[C]{}
%Página \thepage\ de \pageref{LastPage}}% Páginas Pares numeradas a la derecha
%\fancyfoot[LO]{\thepage} % Páginas Impares numeradas a la izquierda

\title{C2_Ay7}

\begin{document}

\begin{center}
Cálculo II\\
Ayudantía Nº7\\
Primer Semestre 2017
\end{center}

\begin{enumerate}
\item Determine si las siguientes afirmaciones son verdaderas o falsas y justifique su respuesta:

\begin{enumerate}[label=\emph{\alph*)}]

\item El valor de $\displaystyle \int_0^\pi f(x) \; dx$, para $f(x)= \left\{ \begin{array}{lcc}
             \sin x &   si  & 0 \leq x < \frac{\pi}{2} \\
             \\ \cos x &  si & \frac{\pi}{2} \leq x \leq \pi 
             \end{array}
   \right.$ es $2\pi$\\

\item El área limitada por las curvas $y = \frac{x}{\sqrt{2x^2+1}}$; $x=0$; $x=2$; $y=0$ es $1$.

\item El valor de $c \in R$, tal que el área acotada por $y = x^2+c$; $x=0$; $x=3$; $y=0$ sea 12, es $c=1$ 
\end{enumerate}
\vspace{2mm}
\item Grafique la regi\'{o}n  limitada por las curvas y calcule su \'{a}rea\\
$\dis y=\frac{1}{x}$ \, ;\,  $\dis y=\frac{1}{x^2}$ \, ;\, $x=1$ \, ;\, $x=2$

\item Grafique la regi\'{o}n  limitada por las curvas y calcule su \'{a}rea\\
$\dis y=\e^x$ \, ;\,  $\dis y=\e^{-x}$ \, ;\, $x=-2$ \, ;\, $x=1$ \, ;\,  $\dis y=0$

\item Calcule el \'{a}rea de la regi\'{o}n limitada por la curva
$y=x^4-2x^3+x^2+3$, el eje $X$ y las rectas $x=a$ ; $x=b$ donde  $a$
y $b$ son las abscisas de los puntos m\'{\i}nimos de la funci\'{o}n.


\item Calcule el \'{a}rea de la regi\'{o}n limitada por la par\'{a}bola
$y=x^2-2x+2$, su tangente en el punto $(3,5)$ y el eje $Y$.




\end{enumerate}




\vspace{10mm}
